\documentclass[11pt,a4paper]{article}
\usepackage[utf8]{inputenc}
\usepackage[T1]{fontenc}
\usepackage[spanish,english]{babel}
\usepackage{lmodern}
\usepackage{textcomp}
\usepackage{geometry}
\usepackage{enumitem}
\usepackage{tikz}
\usepackage{hyperref}
\geometry{margin=2.2cm}
\setlist[enumerate]{leftmargin=*,itemsep=0.8em}
\setlist[itemize]{leftmargin=*,nosep}
\title{MATH CHALLENGE 2 -- SA2: RECYCLING\\Banco completo de preguntas y respuestas}
\author{Documento de revisi\'on docente}
\date{18 de noviembre de 2026}
\begin{document}
\maketitle
\noindent\textbf{Nota de uso.} El banco recoge todas las preguntas disponibles en la versi\'on actual de \texttt{app.js}, no solo las 30 que aparecen en una partida. La aplicaci\'on selecciona aleatoriamente 30 preguntas (10 por bloque); opciones y orden pueden variar. En las preguntas de vocabulario se muestran las opciones de referencia antes de la mezcla aleatoria. Las letras de cada pregunta de unscramble aparecen en una mezcla posible; en la aplicaci\'on se barajan en cada partida.
\tableofcontents
\clearpage
\section{Proportionality and Percentages}
\subsection{Vocabulario: opci\'on m\'ultiple}
\begin{enumerate}
\item What does ``direct proportionality'' mean?
\begin{itemize}
\item proporcionalidad directa
\item proporcionalidad inversa
\item razón
\item proporción
\end{itemize}
\textbf{Respuesta:} proporcionalidad directa
\item What does ``inverse proportionality'' mean?
\begin{itemize}
\item proporcionalidad inversa
\item proporcionalidad directa
\item razón
\item proporción
\end{itemize}
\textbf{Respuesta:} proporcionalidad inversa
\item What does ``ratio'' mean?
\begin{itemize}
\item razón
\item proporcionalidad directa
\item proporcionalidad inversa
\item proporción
\end{itemize}
\textbf{Respuesta:} razón
\item What does ``proportion'' mean?
\begin{itemize}
\item proporción
\item proporcionalidad directa
\item proporcionalidad inversa
\item razón
\end{itemize}
\textbf{Respuesta:} proporción
\item What does ``directly proportional distributions'' mean?
\begin{itemize}
\item repartos directamente proporcionales
\item proporcionalidad directa
\item proporcionalidad inversa
\item razón
\end{itemize}
\textbf{Respuesta:} repartos directamente proporcionales
\item What does ``inversely proportional distributions'' mean?
\begin{itemize}
\item repartos inversamente proporcionales
\item proporcionalidad directa
\item proporcionalidad inversa
\item razón
\end{itemize}
\textbf{Respuesta:} repartos inversamente proporcionales
\item What does ``percentage'' mean?
\begin{itemize}
\item porcentaje
\item proporcionalidad directa
\item proporcionalidad inversa
\item razón
\end{itemize}
\textbf{Respuesta:} porcentaje
\item What does ``percent'' mean?
\begin{itemize}
\item por ciento
\item proporcionalidad directa
\item proporcionalidad inversa
\item razón
\end{itemize}
\textbf{Respuesta:} por ciento
\item What does ``percentage increases'' mean?
\begin{itemize}
\item aumentos porcentuales
\item proporcionalidad directa
\item proporcionalidad inversa
\item razón
\end{itemize}
\textbf{Respuesta:} aumentos porcentuales
\item What does ``percentage decreases'' mean?
\begin{itemize}
\item disminuciones porcentuales
\item proporcionalidad directa
\item proporcionalidad inversa
\item razón
\end{itemize}
\textbf{Respuesta:} disminuciones porcentuales
\item What does ``discount'' mean?
\begin{itemize}
\item descuento
\item proporcionalidad directa
\item proporcionalidad inversa
\item razón
\end{itemize}
\textbf{Respuesta:} descuento
\item What does ``tax'' mean?
\begin{itemize}
\item impuesto
\item proporcionalidad directa
\item proporcionalidad inversa
\item razón
\end{itemize}
\textbf{Respuesta:} impuesto
\item What does ``surcharge'' mean?
\begin{itemize}
\item recargo
\item proporcionalidad directa
\item proporcionalidad inversa
\item razón
\end{itemize}
\textbf{Respuesta:} recargo
\end{enumerate}
\subsection{Unscramble}
\begin{enumerate}
\item Unscramble this mathematical word.\\Letras: o -- i -- t -- a -- r\\Pista: razón
\textbf{Respuesta:} ratio
\item Unscramble this mathematical word.\\Letras: n -- o -- i -- t -- r -- o -- p -- o -- r -- p\\Pista: proporción
\textbf{Respuesta:} proportion
\item Unscramble this mathematical word.\\Letras: e -- g -- a -- t -- n -- e -- c -- r -- e -- p\\Pista: porcentaje
\textbf{Respuesta:} percentage
\item Unscramble this mathematical word.\\Letras: t -- n -- e -- c -- r -- e -- p\\Pista: por ciento
\textbf{Respuesta:} percent
\item Unscramble this mathematical word.\\Letras: t -- n -- u -- o -- c -- s -- i -- d\\Pista: descuento
\textbf{Respuesta:} discount
\item Unscramble this mathematical word.\\Letras: x -- a -- t\\Pista: impuesto
\textbf{Respuesta:} tax
\item Unscramble this mathematical word.\\Letras: e -- g -- r -- a -- h -- c -- r -- u -- s\\Pista: recargo
\textbf{Respuesta:} surcharge
\end{enumerate}
\subsection{Matching}
\begin{enumerate}
\item Match every English word with its Spanish meaning.
\begin{itemize}
\item direct proportionality $\longleftrightarrow$ proporcionalidad directa
\item inverse proportionality $\longleftrightarrow$ proporcionalidad inversa
\item ratio $\longleftrightarrow$ razón
\item proportion $\longleftrightarrow$ proporción
\end{itemize}
\item Match every English word with its Spanish meaning.
\begin{itemize}
\item directly proportional distributions $\longleftrightarrow$ repartos directamente proporcionales
\item inversely proportional distributions $\longleftrightarrow$ repartos inversamente proporcionales
\item percentage $\longleftrightarrow$ porcentaje
\item percent $\longleftrightarrow$ por ciento
\end{itemize}
\item Match every English word with its Spanish meaning.
\begin{itemize}
\item percentage increases $\longleftrightarrow$ aumentos porcentuales
\item percentage decreases $\longleftrightarrow$ disminuciones porcentuales
\item discount $\longleftrightarrow$ descuento
\item tax $\longleftrightarrow$ impuesto
\end{itemize}
\end{enumerate}
\subsection{Fill in the blanks}
\begin{enumerate}
\item A comparison between two quantities is a \_\_\_\_\_\_.
\textbf{Respuesta(s) aceptada(s):} ratio
\item An equality between two ratios is a \_\_\_\_\_\_.
\textbf{Respuesta(s) aceptada(s):} proportion
\item A price reduction is a \_\_\_\_\_\_.
\textbf{Respuesta(s) aceptada(s):} discount
\item An extra charge added to a price is a \_\_\_\_\_\_.
\textbf{Respuesta(s) aceptada(s):} surcharge
\end{enumerate}
\subsection{Interpretaci\'on matem\'atica}
\begin{enumerate}
\item A 3 kg bag costs \texteuro{}6. At the same rate, how much do 5 kg cost?
\begin{itemize}
\item \texteuro{}8
\item \texteuro{}9
\item \texteuro{}10
\item \texteuro{}12
\end{itemize}
\textbf{Respuesta:} \texteuro{}10
\item Share 24 items in the ratio 1:3. How many items are in the smaller share?
\begin{itemize}
\item 4
\item 6
\item 8
\item 18
\end{itemize}
\textbf{Respuesta:} 6
\item A 20\% discount is applied to \texteuro{}50. What is the new price?
\begin{itemize}
\item \texteuro{}10
\item \texteuro{}30
\item \texteuro{}40
\item \texteuro{}60
\end{itemize}
\textbf{Respuesta:} \texteuro{}40
\item A \texteuro{}40 item has 10\% tax added. What is the final price?
\begin{itemize}
\item \texteuro{}4
\item \texteuro{}36
\item \texteuro{}44
\item \texteuro{}50
\end{itemize}
\textbf{Respuesta:} \texteuro{}44
\end{enumerate}
\subsection{True or false}
\begin{enumerate}
\item In direct proportionality, multiplying one quantity by 2 multiplies the other by 2.
\begin{itemize}
\item TRUE
\item FALSE
\end{itemize}
\textbf{Respuesta:} TRUE
\item A 15\% surcharge reduces the price.
\begin{itemize}
\item TRUE
\item FALSE
\end{itemize}
\textbf{Respuesta:} FALSE
\end{enumerate}
\subsection{Detecci\'on de errores}
\begin{enumerate}
\item \textbf{Enunciado:} \textquotedblleft With inverse proportionality, doubling one quantity also doubles the other.\textquotedblright
\textbf{Respuesta:} Incorrecto.
\textbf{Correcci\'on:} With inverse proportionality, doubling one quantity halves the other.
\item \textbf{Enunciado:} \textquotedblleft A 10\% discount lowers the original price.\textquotedblright
\textbf{Respuesta:} Correcto; no contiene error.
\textbf{Correcci\'on:} A 10\% discount lowers the original price.
\end{enumerate}
\section{Decimal Numbers}
\subsection{Vocabulario: opci\'on m\'ultiple}
\begin{enumerate}
\item What does ``decimal number'' mean?
\begin{itemize}
\item número decimal
\item décimas
\item centésimas
\item milésimas
\end{itemize}
\textbf{Respuesta:} número decimal
\item What does ``tenths'' mean?
\begin{itemize}
\item décimas
\item número decimal
\item centésimas
\item milésimas
\end{itemize}
\textbf{Respuesta:} décimas
\item What does ``hundredths'' mean?
\begin{itemize}
\item centésimas
\item número decimal
\item décimas
\item milésimas
\end{itemize}
\textbf{Respuesta:} centésimas
\item What does ``thousandths'' mean?
\begin{itemize}
\item milésimas
\item número decimal
\item décimas
\item centésimas
\end{itemize}
\textbf{Respuesta:} milésimas
\item What does ``ten thousandths'' mean?
\begin{itemize}
\item diezmilésimas
\item número decimal
\item décimas
\item centésimas
\end{itemize}
\textbf{Respuesta:} diezmilésimas
\item What does ``hundred thousandths'' mean?
\begin{itemize}
\item cienmilésimas
\item número decimal
\item décimas
\item centésimas
\end{itemize}
\textbf{Respuesta:} cienmilésimas
\item What does ``millionths'' mean?
\begin{itemize}
\item millonésimas
\item número decimal
\item décimas
\item centésimas
\end{itemize}
\textbf{Respuesta:} millonésimas
\item What does ``terminating decimal'' mean?
\begin{itemize}
\item decimal exacto
\item número decimal
\item décimas
\item centésimas
\end{itemize}
\textbf{Respuesta:} decimal exacto
\item What does ``recurring decimal'' mean?
\begin{itemize}
\item decimal periódico
\item número decimal
\item décimas
\item centésimas
\end{itemize}
\textbf{Respuesta:} decimal periódico
\item What does ``pure recurring decimal'' mean?
\begin{itemize}
\item decimal periódico puro
\item número decimal
\item décimas
\item centésimas
\end{itemize}
\textbf{Respuesta:} decimal periódico puro
\item What does ``mixed recurring decimal'' mean?
\begin{itemize}
\item decimal periódico mixto
\item número decimal
\item décimas
\item centésimas
\end{itemize}
\textbf{Respuesta:} decimal periódico mixto
\item What does ``period / repetend'' mean?
\begin{itemize}
\item periodo
\item número decimal
\item décimas
\item centésimas
\end{itemize}
\textbf{Respuesta:} periodo
\item What does ``non-recurring decimal part'' mean?
\begin{itemize}
\item anteperiodo
\item número decimal
\item décimas
\item centésimas
\end{itemize}
\textbf{Respuesta:} anteperiodo
\item What does ``irrational number'' mean?
\begin{itemize}
\item número irracional
\item número decimal
\item décimas
\item centésimas
\end{itemize}
\textbf{Respuesta:} número irracional
\item What does ``rounding'' mean?
\begin{itemize}
\item redondeo
\item número decimal
\item décimas
\item centésimas
\end{itemize}
\textbf{Respuesta:} redondeo
\item What does ``approximate by rounding to the nearest tenth'' mean?
\begin{itemize}
\item aproxima por redondeo a las décimas
\item número decimal
\item décimas
\item centésimas
\end{itemize}
\textbf{Respuesta:} aproxima por redondeo a las décimas
\item What does ``truncation'' mean?
\begin{itemize}
\item truncamiento
\item número decimal
\item décimas
\item centésimas
\end{itemize}
\textbf{Respuesta:} truncamiento
\item What does ``approximate by truncation to the nearest hundredth'' mean?
\begin{itemize}
\item aproxima por truncamiento a las centésimas
\item número decimal
\item décimas
\item centésimas
\end{itemize}
\textbf{Respuesta:} aproxima por truncamiento a las centésimas
\end{enumerate}
\subsection{Unscramble}
\begin{enumerate}
\item Unscramble this mathematical word.\\Letras: s -- h -- t -- n -- e -- t\\Pista: décimas
\textbf{Respuesta:} tenths
\item Unscramble this mathematical word.\\Letras: s -- h -- t -- d -- e -- r -- d -- n -- u -- h\\Pista: centésimas
\textbf{Respuesta:} hundredths
\item Unscramble this mathematical word.\\Letras: s -- h -- t -- d -- n -- a -- s -- u -- o -- h -- t\\Pista: milésimas
\textbf{Respuesta:} thousandths
\item Unscramble this mathematical word.\\Letras: s -- h -- t -- n -- o -- i -- l -- l -- i -- m\\Pista: millonésimas
\textbf{Respuesta:} millionths
\item Unscramble this mathematical word.\\Letras: g -- n -- i -- d -- n -- u -- o -- r\\Pista: redondeo
\textbf{Respuesta:} rounding
\item Unscramble this mathematical word.\\Letras: n -- o -- i -- t -- a -- c -- n -- u -- r -- t\\Pista: truncamiento
\textbf{Respuesta:} truncation
\end{enumerate}
\subsection{Matching}
\begin{enumerate}
\item Match every English word with its Spanish meaning.
\begin{itemize}
\item decimal number $\longleftrightarrow$ número decimal
\item tenths $\longleftrightarrow$ décimas
\item hundredths $\longleftrightarrow$ centésimas
\item thousandths $\longleftrightarrow$ milésimas
\end{itemize}
\item Match every English word with its Spanish meaning.
\begin{itemize}
\item ten thousandths $\longleftrightarrow$ diezmilésimas
\item hundred thousandths $\longleftrightarrow$ cienmilésimas
\item millionths $\longleftrightarrow$ millonésimas
\item terminating decimal $\longleftrightarrow$ decimal exacto
\end{itemize}
\item Match every English word with its Spanish meaning.
\begin{itemize}
\item recurring decimal $\longleftrightarrow$ decimal periódico
\item pure recurring decimal $\longleftrightarrow$ decimal periódico puro
\item mixed recurring decimal $\longleftrightarrow$ decimal periódico mixto
\item period / repetend $\longleftrightarrow$ periodo
\end{itemize}
\item Match every English word with its Spanish meaning.
\begin{itemize}
\item non-recurring decimal part $\longleftrightarrow$ anteperiodo
\item irrational number $\longleftrightarrow$ número irracional
\item rounding $\longleftrightarrow$ redondeo
\item approximate by rounding to the nearest tenth $\longleftrightarrow$ aproxima por redondeo a las décimas
\end{itemize}
\end{enumerate}
\subsection{Fill in the blanks}
\begin{enumerate}
\item A decimal with a finite number of decimal places is a \_\_\_\_\_\_.
\textbf{Respuesta(s) aceptada(s):} terminating decimal
\item The repeating digits in a recurring decimal form its \_\_\_\_\_\_.
\textbf{Respuesta(s) aceptada(s):} period; repetend; period / repetend
\item The non-repeating decimal part before a repetend is the \_\_\_\_\_\_.
\textbf{Respuesta(s) aceptada(s):} non-recurring decimal part
\item Dropping later digits without rounding is called \_\_\_\_\_\_.
\textbf{Respuesta(s) aceptada(s):} truncation
\end{enumerate}
\subsection{Interpretaci\'on matem\'atica}
\begin{enumerate}
\item In 7.036, what is the place value of the digit 3?
\begin{itemize}
\item tenths
\item hundredths
\item thousandths
\item millionths
\end{itemize}
\textbf{Respuesta:} hundredths
\item Approximate 4.267 by rounding to the nearest tenth.
\begin{itemize}
\item 4.2
\item 4.26
\item 4.3
\item 4.27
\end{itemize}
\textbf{Respuesta:} 4.3
\item Approximate 5.786 by truncation to the nearest hundredth.
\begin{itemize}
\item 5.7
\item 5.78
\item 5.79
\item 5.80
\end{itemize}
\textbf{Respuesta:} 5.78
\item In 0.363636\ldots , which description fits the decimal?
\begin{itemize}
\item terminating decimal
\item pure recurring decimal
\item mixed recurring decimal
\item irrational number
\end{itemize}
\textbf{Respuesta:} pure recurring decimal
\end{enumerate}
\subsection{True or false}
\begin{enumerate}
\item 0.25 is a terminating decimal.
\begin{itemize}
\item TRUE
\item FALSE
\end{itemize}
\textbf{Respuesta:} TRUE
\item An irrational number has a terminating decimal expansion.
\begin{itemize}
\item TRUE
\item FALSE
\end{itemize}
\textbf{Respuesta:} FALSE
\end{enumerate}
\subsection{Detecci\'on de errores}
\begin{enumerate}
\item \textbf{Enunciado:} \textquotedblleft In 0.444\ldots , the digit 4 is a non-recurring decimal part.\textquotedblright
\textbf{Respuesta:} Incorrecto.
\textbf{Correcci\'on:} The digit 4 is the period / repetend.
\item \textbf{Enunciado:} \textquotedblleft Rounding 3.46 to the nearest tenth gives 3.5.\textquotedblright
\textbf{Respuesta:} Correcto; no contiene error.
\textbf{Correcci\'on:} Rounding 3.46 to the nearest tenth gives 3.5.
\end{enumerate}
\section{Geometric Solids}
\subsection{Vocabulario: opci\'on m\'ultiple}
\begin{enumerate}
\item What does ``geometric solid'' mean?
\begin{itemize}
\item cuerpo geométrico
\item poliedro
\item prisma
\item ortoedro
\end{itemize}
\textbf{Respuesta:} cuerpo geométrico
\item What does ``polyhedron'' mean?
\begin{itemize}
\item poliedro
\item cuerpo geométrico
\item prisma
\item ortoedro
\end{itemize}
\textbf{Respuesta:} poliedro
\item What does ``prism'' mean?
\begin{itemize}
\item prisma
\item cuerpo geométrico
\item poliedro
\item ortoedro
\end{itemize}
\textbf{Respuesta:} prisma
\item What does ``orthohedron'' mean?
\begin{itemize}
\item ortoedro
\item cuerpo geométrico
\item poliedro
\item prisma
\end{itemize}
\textbf{Respuesta:} ortoedro
\item What does ``lateral faces'' mean?
\begin{itemize}
\item caras laterales
\item cuerpo geométrico
\item poliedro
\item prisma
\end{itemize}
\textbf{Respuesta:} caras laterales
\item What does ``bases'' mean?
\begin{itemize}
\item bases
\item cuerpo geométrico
\item poliedro
\item prisma
\end{itemize}
\textbf{Respuesta:} bases
\item What does ``height'' mean?
\begin{itemize}
\item altura
\item cuerpo geométrico
\item poliedro
\item prisma
\end{itemize}
\textbf{Respuesta:} altura
\item What does ``triangular prism'' mean?
\begin{itemize}
\item prisma triangular
\item cuerpo geométrico
\item poliedro
\item prisma
\end{itemize}
\textbf{Respuesta:} prisma triangular
\item What does ``quadrangular prism'' mean?
\begin{itemize}
\item prisma cuadrangular
\item cuerpo geométrico
\item poliedro
\item prisma
\end{itemize}
\textbf{Respuesta:} prisma cuadrangular
\item What does ``pentagonal prism'' mean?
\begin{itemize}
\item prisma pentagonal
\item cuerpo geométrico
\item poliedro
\item prisma
\end{itemize}
\textbf{Respuesta:} prisma pentagonal
\item What does ``hexagonal prism'' mean?
\begin{itemize}
\item prisma hexagonal
\item cuerpo geométrico
\item poliedro
\item prisma
\end{itemize}
\textbf{Respuesta:} prisma hexagonal
\item What does ``cube'' mean?
\begin{itemize}
\item cubo
\item cuerpo geométrico
\item poliedro
\item prisma
\end{itemize}
\textbf{Respuesta:} cubo
\item What does ``pyramid'' mean?
\begin{itemize}
\item pirámide
\item cuerpo geométrico
\item poliedro
\item prisma
\end{itemize}
\textbf{Respuesta:} pirámide
\item What does ``vertex'' mean?
\begin{itemize}
\item vértice
\item cuerpo geométrico
\item poliedro
\item prisma
\end{itemize}
\textbf{Respuesta:} vértice
\item What does ``apothem of the pyramid'' mean?
\begin{itemize}
\item apotema de la pirámide
\item cuerpo geométrico
\item poliedro
\item prisma
\end{itemize}
\textbf{Respuesta:} apotema de la pirámide
\item What does ``apothem of the base'' mean?
\begin{itemize}
\item apotema de la base
\item cuerpo geométrico
\item poliedro
\item prisma
\end{itemize}
\textbf{Respuesta:} apotema de la base
\item What does ``solid of revolution'' mean?
\begin{itemize}
\item cuerpo de revolución
\item cuerpo geométrico
\item poliedro
\item prisma
\end{itemize}
\textbf{Respuesta:} cuerpo de revolución
\item What does ``cylinder'' mean?
\begin{itemize}
\item cilindro
\item cuerpo geométrico
\item poliedro
\item prisma
\end{itemize}
\textbf{Respuesta:} cilindro
\item What does ``cone'' mean?
\begin{itemize}
\item cono
\item cuerpo geométrico
\item poliedro
\item prisma
\end{itemize}
\textbf{Respuesta:} cono
\item What does ``axis of revolution'' mean?
\begin{itemize}
\item eje de revolución
\item cuerpo geométrico
\item poliedro
\item prisma
\end{itemize}
\textbf{Respuesta:} eje de revolución
\item What does ``generatrix'' mean?
\begin{itemize}
\item generatriz
\item cuerpo geométrico
\item poliedro
\item prisma
\end{itemize}
\textbf{Respuesta:} generatriz
\item What does ``radius'' mean?
\begin{itemize}
\item radio
\item cuerpo geométrico
\item poliedro
\item prisma
\end{itemize}
\textbf{Respuesta:} radio
\item What does ``area'' mean?
\begin{itemize}
\item área
\item cuerpo geométrico
\item poliedro
\item prisma
\end{itemize}
\textbf{Respuesta:} área
\item What does ``volume'' mean?
\begin{itemize}
\item volumen
\item cuerpo geométrico
\item poliedro
\item prisma
\end{itemize}
\textbf{Respuesta:} volumen
\item What does ``lateral surface area'' mean?
\begin{itemize}
\item área lateral
\item cuerpo geométrico
\item poliedro
\item prisma
\end{itemize}
\textbf{Respuesta:} área lateral
\item What does ``area of the base'' mean?
\begin{itemize}
\item área de la base
\item cuerpo geométrico
\item poliedro
\item prisma
\end{itemize}
\textbf{Respuesta:} área de la base
\end{enumerate}
\subsection{Unscramble}
\begin{enumerate}
\item Unscramble this mathematical word.\\Letras: n -- o -- r -- d -- e -- h -- y -- l -- o -- p\\Pista: poliedro
\textbf{Respuesta:} polyhedron
\item Unscramble this mathematical word.\\Letras: m -- s -- i -- r -- p\\Pista: prisma
\textbf{Respuesta:} prism
\item Unscramble this mathematical word.\\Letras: n -- o -- r -- d -- e -- h -- o -- h -- t -- r -- o\\Pista: ortoedro
\textbf{Respuesta:} orthohedron
\item Unscramble this mathematical word.\\Letras: s -- e -- s -- a -- b\\Pista: bases
\textbf{Respuesta:} bases
\item Unscramble this mathematical word.\\Letras: t -- h -- g -- i -- e -- h\\Pista: altura
\textbf{Respuesta:} height
\item Unscramble this mathematical word.\\Letras: e -- b -- u -- c\\Pista: cubo
\textbf{Respuesta:} cube
\item Unscramble this mathematical word.\\Letras: d -- i -- m -- a -- r -- y -- p\\Pista: pirámide
\textbf{Respuesta:} pyramid
\item Unscramble this mathematical word.\\Letras: x -- e -- t -- r -- e -- v\\Pista: vértice
\textbf{Respuesta:} vertex
\item Unscramble this mathematical word.\\Letras: r -- e -- d -- n -- i -- l -- y -- c\\Pista: cilindro
\textbf{Respuesta:} cylinder
\item Unscramble this mathematical word.\\Letras: e -- n -- o -- c\\Pista: cono
\textbf{Respuesta:} cone
\item Unscramble this mathematical word.\\Letras: x -- i -- r -- t -- a -- r -- e -- n -- e -- g\\Pista: generatriz
\textbf{Respuesta:} generatrix
\item Unscramble this mathematical word.\\Letras: s -- u -- i -- d -- a -- r\\Pista: radio
\textbf{Respuesta:} radius
\item Unscramble this mathematical word.\\Letras: a -- e -- r -- a\\Pista: área
\textbf{Respuesta:} area
\item Unscramble this mathematical word.\\Letras: e -- m -- u -- l -- o -- v\\Pista: volumen
\textbf{Respuesta:} volume
\end{enumerate}
\subsection{Matching}
\begin{enumerate}
\item Match every English word with its Spanish meaning.
\begin{itemize}
\item geometric solid $\longleftrightarrow$ cuerpo geométrico
\item polyhedron $\longleftrightarrow$ poliedro
\item prism $\longleftrightarrow$ prisma
\item orthohedron $\longleftrightarrow$ ortoedro
\end{itemize}
\item Match every English word with its Spanish meaning.
\begin{itemize}
\item lateral faces $\longleftrightarrow$ caras laterales
\item bases $\longleftrightarrow$ bases
\item height $\longleftrightarrow$ altura
\item triangular prism $\longleftrightarrow$ prisma triangular
\end{itemize}
\item Match every English word with its Spanish meaning.
\begin{itemize}
\item quadrangular prism $\longleftrightarrow$ prisma cuadrangular
\item pentagonal prism $\longleftrightarrow$ prisma pentagonal
\item hexagonal prism $\longleftrightarrow$ prisma hexagonal
\item cube $\longleftrightarrow$ cubo
\end{itemize}
\item Match every English word with its Spanish meaning.
\begin{itemize}
\item pyramid $\longleftrightarrow$ pirámide
\item vertex $\longleftrightarrow$ vértice
\item apothem of the pyramid $\longleftrightarrow$ apotema de la pirámide
\item apothem of the base $\longleftrightarrow$ apotema de la base
\end{itemize}
\item Match every English word with its Spanish meaning.
\begin{itemize}
\item solid of revolution $\longleftrightarrow$ cuerpo de revolución
\item cylinder $\longleftrightarrow$ cilindro
\item cone $\longleftrightarrow$ cono
\item axis of revolution $\longleftrightarrow$ eje de revolución
\end{itemize}
\item Match every English word with its Spanish meaning.
\begin{itemize}
\item generatrix $\longleftrightarrow$ generatriz
\item radius $\longleftrightarrow$ radio
\item area $\longleftrightarrow$ área
\item volume $\longleftrightarrow$ volumen
\end{itemize}
\end{enumerate}
\subsection{Fill in the blanks}
\begin{enumerate}
\item The faces that are not the bases of a prism are its \_\_\_\_\_\_.
\textbf{Respuesta(s) aceptada(s):} lateral faces
\item The perpendicular distance between the bases is the \_\_\_\_\_\_.
\textbf{Respuesta(s) aceptada(s):} height
\item The area of the side faces together is the \_\_\_\_\_\_.
\textbf{Respuesta(s) aceptada(s):} lateral surface area
\item The measure of one base is the \_\_\_\_\_\_.
\textbf{Respuesta(s) aceptada(s):} area of the base
\end{enumerate}
\subsection{Interpretaci\'on matem\'atica}
\begin{enumerate}
\item A prism has a triangular base. What is its name?
\begin{itemize}
\item triangular prism
\item quadrangular prism
\item pentagonal prism
\item hexagonal prism
\end{itemize}
\textbf{Respuesta:} triangular prism
\item Which measurement describes the distance from the centre of a circular base to its edge?
\begin{itemize}
\item height
\item generatrix
\item radius
\item vertex
\end{itemize}
\textbf{Respuesta:} radius
\end{enumerate}
\subsection{True or false}
\begin{enumerate}
\item A cube is a polyhedron.
\begin{itemize}
\item TRUE
\item FALSE
\end{itemize}
\textbf{Respuesta:} TRUE
\item A cone has two parallel bases.
\begin{itemize}
\item TRUE
\item FALSE
\end{itemize}
\textbf{Respuesta:} FALSE
\end{enumerate}
\subsection{Detecci\'on de errores}
\begin{enumerate}
\item \textbf{Enunciado:} \textquotedblleft A cube has six square faces.\textquotedblright
\textbf{Respuesta:} Correcto; no contiene error.
\textbf{Correcci\'on:} A cube has six square faces.
\item \textbf{Enunciado:} \textquotedblleft A cylinder is a polyhedron because it has curved surfaces.\textquotedblright
\textbf{Respuesta:} Incorrecto.
\textbf{Correcci\'on:} A polyhedron has only flat polygonal faces; a cylinder has a curved surface.
\end{enumerate}
\subsection{Reconocimiento visual}
\begin{enumerate}
\item Which geometric solid is shown?
\begin{center}
\begin{tikzpicture}[scale=0.55,line width=0.8pt]
\draw (0,0)--(2,0)--(2,2)--(0,2)--cycle;
\draw (0.8,0.8)--(2.8,0.8)--(2.8,2.8)--(0.8,2.8)--cycle;
\draw (0,0)--(0.8,0.8) (2,0)--(2.8,0.8) (2,2)--(2.8,2.8) (0,2)--(0.8,2.8);
\end{tikzpicture}
\end{center}
\begin{itemize}
\item cube
\item triangular prism
\item pyramid
\item cylinder
\item cone
\end{itemize}
\textbf{Respuesta:} cube
\item Which geometric solid is shown?
\begin{center}
\begin{tikzpicture}[scale=0.55,line width=0.8pt]
\draw (0,0)--(1,1.8)--(2,0)--cycle;
\draw (1,0.5)--(2,2.3)--(3,0.5)--cycle;
\draw (0,0)--(1,0.5) (1,1.8)--(2,2.3) (2,0)--(3,0.5);
\end{tikzpicture}
\end{center}
\begin{itemize}
\item cube
\item triangular prism
\item pyramid
\item cylinder
\item cone
\end{itemize}
\textbf{Respuesta:} triangular prism
\item Which geometric solid is shown?
\begin{center}
\begin{tikzpicture}[scale=0.55,line width=0.8pt]
\draw (0,0)--(1,0.5)--(2,0)--(1,-0.5)--cycle;
\draw (1,2.2)--(0,0) (1,2.2)--(1,0.5) (1,2.2)--(2,0) (1,2.2)--(1,-0.5);
\end{tikzpicture}
\end{center}
\begin{itemize}
\item cube
\item triangular prism
\item pyramid
\item cylinder
\item cone
\end{itemize}
\textbf{Respuesta:} pyramid
\item Which geometric solid is shown?
\begin{center}
\begin{tikzpicture}[scale=0.55,line width=0.8pt]
\draw (0,0) ellipse (1 and 0.3);
\draw (0,2) ellipse (1 and 0.3);
\draw (-1,0)--(-1,2) (1,0)--(1,2);
\end{tikzpicture}
\end{center}
\begin{itemize}
\item cube
\item triangular prism
\item pyramid
\item cylinder
\item cone
\end{itemize}
\textbf{Respuesta:} cylinder
\item Which geometric solid is shown?
\begin{center}
\begin{tikzpicture}[scale=0.55,line width=0.8pt]
\draw (0,2)--(-1,0) (0,2)--(1,0);
\draw (0,0) ellipse (1 and 0.3);
\end{tikzpicture}
\end{center}
\begin{itemize}
\item cube
\item triangular prism
\item pyramid
\item cylinder
\item cone
\end{itemize}
\textbf{Respuesta:} cone
\end{enumerate}
\end{document}
